\chapter*{Abstract}
\addcontentsline{toc}{chapter}{\MakeUppercase{Abstract}}

\small Abstract is a concise summary of the entire final project, containing the background, objectives, research methods, results, and conclusions. Its purpose is to provide readers with an overview of the study, enabling them to quickly understand the research problem, the approach employed, and the outcomes achieved without reading the entire document.

An abstract is typically written in one to three paragraphs, depending on the applicable guidelines. It should be concise, clear, and informative, and should not contain citations, tables, figures, or equations.

In general, the abstract may be organized as follows:

\begin{itemize}[topsep=0pt,itemsep=0pt,parsep=0pt,leftmargin=*]
\item The first paragraph describes the research background, objectives, and scope or limitations of the study.
\item The second paragraph explains the methods, approaches, algorithms, or procedures used to address the research problem.
\item The third paragraph presents the main findings, the level of achievement attained, and the conclusions drawn. Whenever possible, the results should be presented quantitatively using relevant performance indicators or evaluation metrics.
\end{itemize}

An abstract should be able to answer the following questions:

\begin{enumerate}[topsep=0pt,itemsep=0pt,parsep=0pt,leftmargin=*]
\item What problem is being investigated?
\item Why is the problem important?
\item How was the research conducted?
\item What results were obtained?
\item What are the main conclusions of the study?
\end{enumerate}

After the abstract, include several keywords that represent the main topics of the research to facilitate indexing and document retrieval.\\
\vspace{1em}

\noindent \textbf{Keywords:} List two to six keywords related to the topic being discussed.

\normalsize
